\documentclass[10pt,a4paper]{amsart}
\usepackage[margin=19mm]{geometry}
\usepackage{amsmath,amssymb,mathtools}
\usepackage[scr=rsfs,cal=euler]{mathalfa}
\usepackage{xfrac}
\usepackage[unicode,hidelinks,bookmarks=false]{hyperref}
\newcommand{\ie}{i.e.\ }
\newcommand{\cf}{cf.\ }
\newcommand{\resp}{respectively\ }
\newcommand{\Subsection}{\subsection}
\providecommand{\nobreakdash}{\nobreak\hbox{-}}
\setlength{\emergencystretch}{2em}
\multlinegap0pt
\usepackage{bm}
\usepackage{bbm}
\usepackage{dsfont}
\usepackage{xspace}
\usepackage{cancel}
\usepackage{mathtools}
\usepackage{amsthm}
\makeatletter
\@ifundefined{theorem}{\newtheorem{theorem}{Theorem}}{}
\@ifundefined{lemma}{\newtheorem{lemma}[theorem]{Lemma}}{}
\makeatother
\title[Young-lattice diagonals and a doubly graded multiple-zeta de]{Young-lattice diagonals and a doubly graded multiple-zeta decomposition of $e^\gamma$}
\author{Ricardo G\'omez-A\'iza}
\date{Source: arXiv:2608.22561v1 (CC BY 4.0)}
\begin{document}
\maketitle

\noindent\emph{Source and licence.} Young-lattice diagonals and a doubly graded multiple-zeta decomposition of $e^\gamma$. Version arXiv:2608.22561v1 (CC BY 4.0), distributed under the Creative Commons Attribution 4.0 International licence, as recorded in the arXiv licence metadata for this exact version and shown on the abstract page https://arxiv.org/abs/2608.22561v1. Full rights notice: licenses/arxiv-2608.22561-CC-BY-4.0.txt.\par\smallskip

\noindent\textit{Editorial scope.} Counted unit: the Theorem whose source label is \texttt{thm:CnuMZV} in the article (source0.tex lines 2217--2245, source label \texttt{thm:CnuMZV}), delivered together with the article's own complete original proof of it (source0.tex lines 2251--2313, one \texttt{proof} environment of 63 lines). No mathematics is written, repaired or re-derived: statement and proof are the article's own text line for line. The proof is detached from the statement in the source: the article states the result at lines 2217--2245 and prints its proof at lines 2251--2313, and the proof environment's own header names the statement, so the association is the article's own. Required context retained uncounted: eq:Hcal-t (lines 1989--1999); eq:geom-reg (lines 2005--2012); lem:regabs (lines 2016--2030). Every definition, display and label of those lines is retained unchanged. Omitted from this package and not redistributed: 1--1988, 2000--2004, 2013--2015, 2031--2216, 2246--2250, 2314--2693. Nothing inside a retained excerpt is omitted or altered: every retained body is delivered exactly as the article's lines are written, so no omission receipt of any kind accompanies this package. Citations: the retained excerpts carry no citation command at all, so no bibliography entry of the article is transcribed and no reference is redistributed. Reference adaptation: none was needed and none was made; every reference inside the delivered unit resolves inside it, so no unresolved reference or citation is produced. Printed slips kept literal: no printed word, symbol, hypothesis or number of the counted statement or of its proof was altered. Layout and class: the article's own class is \texttt{article} with packages including geometry, amsmath, amssymb, amsthm, mathtools, bm, mathrsfs, booktabs, enumitem, hyperref, microtype; the wrapper is the lane's pinned \texttt{amsart} editorial wrapper at 10pt on A4, which restates the theorem-like environments the retained text needs and adds the article's own macro definitions that the retained text needs (none), with extra wrapper packages (bm, bbm, dsfont, xspace, cancel, mathtools). The delivered numbering is the unit's own. Assets: no retained span contains a figure environment, a graphics-inclusion command, a PDF-page inclusion, a table or a TikZ picture; no image file is copied into this package, the package declares no asset dependency and no image file is redistributed with it. Licence: version arXiv:2608.22561v1 of this article is distributed under the Creative Commons Attribution 4.0 International licence, as recorded in the arXiv licence metadata for this exact version and shown on the abstract page https://arxiv.org/abs/2608.22561v1. A full rights notice is delivered at licenses/arxiv-2608.22561-CC-BY-4.0.txt. Characters: every retained excerpt is pure ASCII, so the delivered engine, XeLaTeX with its default Latin Modern text font, reports no missing character.
\begin{equation}\label{eq:Hcal-t}
 \mathcal H_t(u)
 :=
 \prod_{q=1}^\infty
 \left(
 1+
 \sum_{a\geq2}
 \frac{u^{a-1}}
 {a!\,q^{a-1}(q+t)}
 \right).
\end{equation}

\begin{equation}\label{eq:geom-reg}
 \frac{1}{q^{a-1}(q+t)}
 =
 \frac1{q^a}\frac1{1+t/q}
 =
 \sum_{m\geq0}
 \frac{(-t)^m}{q^{a+m}}.
\end{equation}

\begin{lemma}[Absolute convergence with regulator]
\label{lem:regabs}
Fix $0\leq t<1$ and $R>0$.  Then
\begin{equation}\label{eq:regabs}
 \sum_{q\geq1}
 \sum_{a\geq2}
 \sum_{m\geq0}
 \frac{R^{a-1}t^m}
 {a!\,q^{a+m}}
 <\infty.
\end{equation}
Consequently, for $|u|\leq R$, the sums and products obtained by
expanding \eqref{eq:Hcal-t} and \eqref{eq:geom-reg} may be rearranged
absolutely.
\end{lemma}

\begin{theorem}[Partition--MZV correspondence]
\label{thm:CnuMZV}
For every partition
$\nu=(\nu_1,\ldots,\nu_d)\vdash r-1$,
\begin{align}
 C_\nu
 =
 \lim_{t\uparrow1}
 \frac{1}
 {\prod_{j=1}^d(\nu_j+1)!}
 \sum_{\sigma\in\operatorname{Perm}(\nu)}
 \sum_{m_1,\ldots,m_d\geq0}
 &
 (-t)^{m_1+\cdots+m_d}
 \nonumber\\
 {}\times&
 \zeta\!\left(
 \nu_{\sigma(1)}+1+m_1,\ldots,
 \nu_{\sigma(d)}+1+m_d
 \right).
 \label{eq:CnuMZV}
\end{align}
In particular, the length
\[
        d=\ell(\nu)
\]
is the depth of every MZV occurring in the block associated with
$\nu$.
\end{theorem}

\begin{proof}[Proof of Theorem \ref{thm:CnuMZV}]
For $0\leq t<1$, define
\[
 C_\nu(t)
 :=
 \frac{1}
 {\prod_{j=1}^d(\nu_j+1)!}
 \sum_{\sigma\in\operatorname{Perm}(\nu)}
 \sum_{1\leq q_1<\cdots<q_d}
 \prod_{j=1}^d
 \frac{1}
 {q_j^{\nu_{\sigma(j)}}(q_j+t)}.
\]
Applying \eqref{eq:geom-reg} to each factor gives
\begin{align*}
 C_\nu(t)
 =
 \frac{1}
 {\prod_{j=1}^d(\nu_j+1)!}
 \sum_{\sigma\in\operatorname{Perm}(\nu)}
 \sum_{m_1,\ldots,m_d\geq0}
 &
 (-t)^{m_1+\cdots+m_d}\\
 {}\times&
 \zeta\!\left(
 \nu_{\sigma(1)}+1+m_1,\ldots,
 \nu_{\sigma(d)}+1+m_d
 \right),
\end{align*}
where absolute convergence follows from
Lemma~\ref{lem:regabs}.

On the other hand, for $0\leq t\leq1$,
\[
 \frac{1}
 {q_j^{\nu_{\sigma(j)}}(q_j+t)}
 \leq
 \frac{1}
 {q_j^{\nu_{\sigma(j)}+1}}.
\]
Hence
\[
 \prod_{j=1}^d
 \frac{1}
 {q_j^{\nu_{\sigma(j)}}(q_j+t)}
 \leq
 \prod_{j=1}^d
 \frac{1}
 {q_j^{\nu_{\sigma(j)}+1}}.
\]
Since every $\nu_{\sigma(j)}+1\geq2$, the series
\[
 \sum_{1\leq q_1<\cdots<q_d}
 \prod_{j=1}^d
 q_j^{-\nu_{\sigma(j)}-1}
\]
converges absolutely.  Dominated convergence therefore gives
\[
        C_\nu(t)\longrightarrow C_\nu
        \qquad (t\uparrow1),
\]
which proves \eqref{eq:CnuMZV}.
\end{proof}

\end{document}
